\bigskip

\textbf{Affirmation 1 : Fausse}

Étudions les variations de la fonction $f$.

$f$ est dérivable sur $ [0~;~+\infty[ $, en tant que composée de fonctions dérivables sur $\R$.

Pour tout $t$ réel positif, on pose : $u(t) = -0,5t^2 + t+2$, donc on a $u'(t) = -t + 1$.

On a donc, pour tout réel $t$ positif : $f'(t) = 0 - (-t + 1)\e^{-0,5t^2 + t + 2} = (t-1)\e^{-0,5t^2 + t + 2} $

La fonction exp étant à valeurs strictement positives sur $\R$, le signe de $f'(t)$ est le signe de $(t-1)$ : c'est donc strictement négatif sur l'intervalle $ [0~;~1[ $ et strictement positif sur $ ]1~;~+\infty[ $.

$f$ est donc strictement décroissante sur l'intervalle $[ 0~;~1]$ : c'est pourquoi il est faux de dire que la population augmente en permanence : au cours de la première heure, la population décroit, strictement.

\medskip
\textbf{Affirmation 2 : Fausse}

La fonction exponentielle est à valeurs strictement positives, donc :

$\begin{aligned} \forall t \in \R^+, \quad \e^{-0,5t^2 + t + 2} > 0 &\implies -\e^{-0,5t^2 + t + 2}<0\\
&\implies \e^3-\e^{-0,5t^2 + t + 2}<\e^3\\
&\implies f(t) < \e^3\\ \end{aligned}$

Or $\e^3 \approx 20,086$ (à $ 10^{-3} $ près, arrondi par excès), donc la fonction $f$ est majorée par 20,086, ce qui s'interprète en disant que la population restera inférieure à 20,086 milliers d'individus, soit \np{20086} bactéries : la population ne dépassera donc jamais \np{21000}.

\medskip

\textbf{Affirmation 3 : Vraie}

Une population de bactéries de \np{10000} individus à l'instant $t$, c'est équivalent à avoir $f(t) = 10$.

On cherche donc à savoir combien d'antécédents $10$ a par la fonction $f$.

En reprenant l'étude des variations de $f$ : on a établi que $f$ était strictement décroissante sur $[0~;~1]$, puis strictement croissante sur $ [1~;~+\infty[ $.

Étudions la limite de $f$ en $+\infty$.

$\lim\limits_{t \to +\infty} -0,5t^2 + t +2 = \lim\limits_{t \to +\infty} -0,5t^2 \left(1 - \dfrac{2}{t} - \dfrac{4}{t}\right) = -\infty$, par limite de quotient, somme et produit de suites de références.

Par composition, en posant $y = -0,5t^2 + t +2$, on a $\lim\limits_{t \to+\infty} \e^{-0,5t^2 + t +2} = \lim\limits_{y\to -\infty} \e^y = 0$

Finalement, par limite de la somme $\lim\limits_{t \to +\infty} f(t) = \e^3$.

Pour rappel, on a $\e^3 \approx 20,086 > 10$.

Par ailleurs : $f(0) = \e^3 - \e^2 \approx 12,696 >10$ et $f(1) = \e^3 - \e^{2,5} \approx 7,903 < 10$.

On peut donc dire :

\begin{itemize}
\item Sur $ [0~;~1] $, $f$ est une fonction continue (car dérivable), et strictement décroissante, et 10 est une valeur intermédiaire entre $f(0)$ et $f(1)$, donc en vertu du corollaire au théorème des valeurs intermédiaires appliqué aux fonctions strictement monotones, l'équation $f(t)=10$ admet une unique solution sur l'intervalle $ [0~;~1] $.
\item Sur $ ]1~;~+\infty[ $, $f$ est une fonction continue (car dérivable), et strictement croissante, et 10 est une valeur intermédiaire strictement comprise entre $\lim\limits_{t \to 1}f(t)=f(1)$ et $\lim\limits_{t\to +\infty}f(t)$, donc en vertu du corollaire au théorème des valeurs intermédiaires appliqué aux fonctions strictement monotones, l'équation $f(t)=10$ admet une unique solution sur l'intervalle $]1~;~+\infty[$.
\end{itemize}

Finalement, l'équation $f(t) = 10$ admet exactement deux solutions sur $\R^+$ et donc il y a exactement deux instants où la population de bactéries aura un effectif de \np{10000} individus.
